\chapter{失敗と教訓}\label{ch:lessons}

この章には、2 週間の作業で起きた誤り、設計の穴、手順の失敗を、見つけた経緯とともにまとめる。どれも各章に書いてあるが、まとめて読むと、同じ種類の誤りがくり返されていることがわかる。

\section{見つけた誤り}

\begin{longtable}{P{3.2cm}P{5.6cm}P{5.6cm}}
\toprule
\textbf{誤り} & \textbf{中身} & \textbf{見つけた経緯と対処} \\
\midrule
\endhead
\bottomrule
\endfoot
\multicolumn{3}{l}{\textbf{単位}} \\
無リスク金利の単位（元の実装） & 日次のリターンから年率 2\% をそのまま引き、均衡リターンが「毎日 1〜2\% 下がる」値になった & 再現の途中で、均衡リターンの値を見て気づいた。日次の $0.02/252$ に直した \\
$\alpha$ の $\sqrt{21}$（実験 7・9） & 1 か月の見通しを日次の $\sigma$ で作り、$\alpha$ が期間に合った大きさの 4.6 倍になった & 実験の後、利用者の「日次のボラティリティを使っていないか」という指摘で見直した。TE の上限が効いていたので結果は同じだった（X05） \\
\midrule
\multicolumn{3}{l}{\textbf{計算}} \\
最適化の打ち切り（元の実装） & SLSQP が「収束した」と報告して途中で止まり、その位置が入力のわずかな差で変わる & H01 で再現できないことから見つけた。Clarabel で最後まで解く形に直した \\
提出日の取り違え（DB） & 同じ四半期に 2 種類の項目名があるとき、翌年の報告書の提出日を選んでいた（APA） & 実験 C で、決算の古さの最大値（487 日）を調べて見つけた。最初に提出された値を優先する形に直した \\
判断日の一覧（実験 9） & 事前登録は 2016-03 からだったが、X02 から借りた関数が 2017-01 からしか日付を返さなかった & 実行の後に気づいた。主な判定は元の実行で行い、日付を直した実行も報告した（結論は同じ） \\
\midrule
\multicolumn{3}{l}{\textbf{設計}} \\
QLIKE の比較（X02） & lw252 で最適化した比率のリスクを、ほかの推定と QLIKE で比べた。lw252 が不利になるのは仕組み上の結果（最適化の呪い） & 実行の後に気づき、そう明記した。判定は実現の TE で行った \\
修正 1 の規則（実験 11） & 定型文の語を含む行を、行ごと捨てていた。要約が 1 行に詰まった発表では、大事な中身ごと消えた（約 15\% の発表） & 開発の判定が、文章が空になった 21 本で止まったことから見つけた。どの評価よりも前だったので、文ごとに捨てる規則に直し、追記として記録した \\
ロジスティック回帰（実験 11） & 外れ値に引っぱられ、アナリスト予想との差だけの順位にも負けた & 最後の比較で見つけた。順位にした版を探索として加えた \\
\midrule
\multicolumn{3}{l}{\textbf{手順}} \\
RunPod の結果の消失（実験 11） & 結果をダウンロードする前に Pod を止め、ディスクごと消えた & この PC で約 19 時間かけてやり直した \\
撤回された論文の引用 & Kim・Muhn・Nikolaev（2024）の論文を成功例として引用していたが、2025 年 2 月に撤回されていた & 利用者の依頼で文献を調べ直して見つけ、引用したすべての文書を直した（F1） \\
LLM の文献調査の誤り & GPT・Gemini・Deep Research の文献調査に、別の論文の番号、著者の取り違え、要旨にない数値があった & 文献を arXiv・OpenAlex・Crossref と照合するスクリプト（\code{refcheck.py}）で見つけた \\
\end{longtable}

\section{教訓}

\begin{enumerate}
\item \textbf{単位は、式を書いた時点で確かめる}。日次・月次・年率の取り違えは 2 回起きた。均衡リターン、$\alpha$、共分散、無リスク金利の単位を、コードの 1 か所にまとめて書き、テストで確かめるべきだった。
\item \textbf{「収束した」を信じない}。最適化の状態の報告だけでなく、別の解法で同じ問題を解き直して目的関数の値を比べる（H01 で OSQP と比べたように）。
\item \textbf{事前登録は、誤りも見えるようにする}。実験 9 の日付の誤りは、事前登録の文書と実行を比べたから見つかった。計画がなければ、ずれに気づかないか、気づいても黙って直していたかもしれない。
\item \textbf{評価の前の確認が効く}。修正 1 の規則の穴は、株価の反応を一度も見る前に見つかったので、直しても結果を選んだことにならなかった。データの加工の段階で、少数の実例（Abbott の発表文、CenterPoint の発表文）を人が読む価値は大きい。
\item \textbf{仕組みの上で決まる比較をしない}。X02 の QLIKE の比較は、比べる前から答えが決まっていた。比較の設計の段階で「この比較は、ある候補に有利になる仕組みがないか」を問う。
\item \textbf{借りた計算機の結果は、止める前に手元に移す}。
\item \textbf{LLM の挙げた文献は、必ず照合する}。存在しても、中身（数値、結論）が違うことがある。撤回も確かめる。
\item \textbf{スコープを絞る}。追加のデータを集め始めたとき、「横に広げすぎていないか」を見直したことで、実験 11 の問い（何が効いたか、どこを直せるか）に集中できた。
\end{enumerate}

\section{うまくいったこと}

\begin{itemize}
\item \textbf{事前登録と 1 回だけの実行}：実験 2〜11 の判定は、すべて結果を見る前に固定した規則で出した。
\item \textbf{置換の対照}：同じ日の中で点数を入れ替える対照は、分布の仮定に頼らず、日付の効果を保つ。$t$ 検定と置換の両方で判定したので、結果が強い。
\item \textbf{比べる相手を強くする}：gpt-4o-mini の点数を、反応で学習したモデル（TF-IDF、FinBERT の追加学習）と比べた。学習なしの LLM に対して、学習したモデルは有利な条件である。
\item \textbf{先読みのないモデル}：ChronoGPT の版を使い、LLM の結果が記憶でないことを確かめた。
\item \textbf{別のモデルにレビューさせる}：実験 1 の結果、設計、コードを GPT と Gemini にレビューさせ、指摘を 1 つずつ確かめた（付録）。
\end{itemize}
